\documentclass[11pt,a4paper]{article}

\usepackage[utf8]{inputenc}
\usepackage[T1]{fontenc}
\usepackage[spanish]{babel}
\usepackage[a4paper, top=2.3cm, bottom=2.3cm, left=2cm, right=2cm, headheight=15pt]{geometry}
\usepackage{xcolor}
\usepackage{tabularx}
\usepackage{booktabs}
\usepackage{array}
\usepackage{fancyhdr}
\usepackage{lastpage}

\definecolor{AzulPericial}{RGB}{27, 54, 93}
\definecolor{GrisTexto}{RGB}{80, 80, 80}

\pagestyle{fancy}
\fancyhf{}
\renewcommand{\headrulewidth}{0.5pt}
\renewcommand{\footrulewidth}{0.5pt}
\fancyhead[L]{\small\textcolor{GrisTexto}{Informe Técnico Pericial}}
\fancyhead[R]{\small\textcolor{GrisTexto}{Ref.: ITP26-GP08}}
\fancyfoot[L]{\small\textcolor{GrisTexto}{Grupo de Prácticas H --- Peritos Ingenieros en Informática}}
\fancyfoot[R]{\small\textcolor{GrisTexto}{Página \thepage\ de \pageref{LastPage}}}

\setlength{\parindent}{0pt}
\setlength{\parskip}{0.4em}

\begin{document}

\section*{\textcolor{AzulPericial}{Anexo II. Distribución de actuaciones del Equipo Pericial}}

A efectos de trazabilidad interna y transparencia metodológica, se detalla a continuación la distribución de responsabilidades técnicas, custodia de soportes y análisis forenses ejecutados de forma colegiada por cada uno de los integrantes del \textbf{Equipo Pericial Grupo de Prácticas H} en la elaboración del presente dictamen:

\vspace{0.35cm}
\renewcommand{\arraystretch}{1.35}
\begin{tabularx}{\textwidth}{@{} >{\raggedright\arraybackslash}p{3.9cm} >{\raggedright\arraybackslash}p{3.5cm} X @{}}
    \toprule
    \textbf{\textcolor{AzulPericial}{Perito / Colegiado}} & 
    \textbf{\textcolor{AzulPericial}{Área de especialidad}} & 
    \textbf{\textcolor{AzulPericial}{Actuaciones y evidencias asignadas}} \\
    \midrule
    \textbf{Pablo Pastor Aguado} \newline {\small Col. n.\textsuperscript{o} 12} & 
    Arquitectura de sistemas y soportes extraíbles & 
    \textbf{Escenario 3 (Dispositivo USB):} Custodia del soporte físico \texttt{General UDisk}, habilitación de bloqueo contra escritura (\texttt{WriteProtect = 1}), adquisición bit a bit en formato RAW (\texttt{usb\_evidencia.001}) con \textit{FTK Imager} y recuperación de \texttt{Confidencial.pdf} desde la \texttt{\$MFT} mediante \textit{Autopsy} (Sección 7). \\
    \midrule
    \textbf{Hugo Adán de la Fuente} \newline {\small Col. n.\textsuperscript{o} 12} & 
    Redes, protocolos SMTP y correo electrónico & 
    \textbf{Escenario 1 (Correo electrónico):} Custodia del disco virtual \texttt{Maquina\_Corrreo.vdi}, descifrado de volumen BitLocker en solo lectura (\textit{dislocker} en Kali Linux), extracción del mensaje MIME (\texttt{.eml}), análisis de \texttt{envios.log} y reconstrucción cronológica (Sección 5). \\
    \midrule
    \textbf{Santiago González Rodríguez} \newline {\small Col. n.\textsuperscript{o} 12} & 
    Sistemas de archivos NTFS y entornos en la nube & 
    \textbf{Escenario 2 (Dropbox):} Examen forense de la imagen \texttt{W11.dd}, inspección de la entrada MFT \texttt{188015} mediante \textit{Autopsy} y \textit{The Sleuth Kit} (\texttt{istat}), verificación del flujo alternativo \texttt{com.dropbox.attrs} y trazabilidad en \texttt{\$Recycle.Bin} (Sección 6). \\
    \midrule
    \textbf{Álvaro López Primo} \newline {\small Col. n.\textsuperscript{o} 12} & 
    Seguridad de la información y cadena de custodia & 
    \textbf{Infraestructura, metodología y consolidación:} Implementación y auditoría del servidor SMTP/Webmail de pruebas (\texttt{IFOR\_P2}), control de verificación criptográfica (SHA-256), adecuación normativa (ISO/IEC 27037 y UNE 71506) y síntesis de conclusiones (Secciones 3, 4, 8 y 9). \\
    \bottomrule
\end{tabularx}

\end{document}
